\section{The twisted congruence defect}\label{sec:twisted-defect}

Fix a mesh point $x$, the weight set $T$ and the shift $u_0$ of
\S\ref{sec:theta-chain}. At a twisted place the projection of an
integral form onto the lifted eigenspaces can have a denominator.
The period comparison needs the length of this projection lattice
modulo its integral part. Put
\[
 y=x^\varepsilon,\qquad N_c=\frac{N_2kby^2}{n^2},\qquad
 B_0=fmby,\qquad m(B_0+N_c)\le\mathcal B .
\]
All words and all $m$ lifts below belong to this shift. Write
$d_E=[E:\mathbb Q]$ and $\Delta_E=|\operatorname{disc}E|$.

\subsection{Projection and trimming}

Let $w$ be a place of the coefficient field $F$ of residue
characteristic $l$. Use $\ell_w=\log N\mathfrak p_w/[F:\mathbb Q]$
and $e_w=v_w(l)$ as in \S\ref{sec:inert-gain}. At a twisted place,
localize the integral Hecke lattice of lift $j$ at a maximal ideal of
its Hecke algebra. Write $M$ for this lattice over a complete splitting
discrete valuation ring $O$, $K$ for its fraction field, $v$ for its
integer valuation, and $e_O=v(l)$. In the split semisimple space
$KM=\bigoplus_{a\in X}V_a$, let $G\subseteq X$ consist of the
eigensystems that occur in the lift. For any subset $S$ write
$V_S=\bigoplus_{a\in S}V_a$ and let $\pi_S$ be its spectral projection.
Define
\[
 Q(M,G)=\frac{\pi_G M}{M\cap V_G},\qquad
 \eta(M,G)=\frac{\operatorname{length}_O Q(M,G)}{e_O}.
\]
Both lattices span $V_G$. Flat extension to a larger splitting ring
multiplies the length and $e_O$ by its relative ramification index.
At $w$ sum $\eta(M,G)$ over the local components to obtain $\eta_j(w)$.
Use zero at untwisted places. On the finite twisted support put
\begin{equation}\label{eq:twisted-sum}
 H_j=\sum_w\ell_we_w\eta_j(w),\qquad
 H_{j,l}=\sum_{w\mid l}\ell_we_w\eta_j(w).
\end{equation}
These are the normalizations of \citep[Definition~17.1]{adriazola-grh}.

\begin{lemma}\label{lem:twisted-trim}
For any $S\subseteq X$, projection induces a surjection
$Q(M,G)\to Q(\pi_S M,G\cap S)$. If $G\subseteq S$ and an integral
endomorphism acts as $a$ on $V_G$ and as zero on $V_{X\setminus S}$,
where $a$ is a positive integer, then
\[
 0\le\eta(M,G)-\eta(\pi_S M,G)\le v_l(a)\dim_K V_G .
\]
\end{lemma}

\begin{proof}
The map $\pi_{G\cap S}$ sends $\pi_G M$ onto
$\pi_{G\cap S}\pi_S M$ and sends $M\cap V_G$ into $\pi_S M\cap V_{G\cap S}$.
If $G\subseteq S$, the two numerator lattices agree. For
$z\in\pi_S M\cap V_G$ choose $z+b\in M$ with $b\in V_{X\setminus S}$.
The endomorphism sends $z+b$ to $az\in M\cap V_G$. The difference
of the lengths is the length of $(\pi_S M\cap V_G)/(M\cap V_G)$,
at most $v(a)\dim_K V_G$. Divide by $e_O$.
\end{proof}

Write $d_j^0=\dim\operatorname{im}\theta_j^*$ for the dimension of
the lifted space of standard type, where $\theta_j$ is the standard
theta map to the level-$j$ row space. We use the following trim on
the actual lattices \citep[Propositions~17.32--17.35]{adriazola-grh}
\citep{kimi-defect,kimi-generator}.

\begin{hypothesis}[\openinput]\label{hyp:twisted-lattices}
The integral lattices for the period maps have the components, good
eigenspaces and support in \eqref{eq:twisted-sum}. For every $j,w$,
the sum of $\dim V_G$ over its components is at most $d_j^0$.
There are trims retaining $G$ and endomorphisms as in
Lemma~\ref{lem:twisted-trim}, with a common positive integer $a_0$
at the mesh point and $\log a_0\ll_\varepsilon x^{C_\varepsilon}$.
All size constants in this section are uniform on the chosen windows.
\end{hypothesis}

Denote the trimmed sums by $H_j^{\rm cut},H_{j,l}^{\rm cut}$, with
the places and weights of \eqref{eq:twisted-sum}. Since
$\sum_{w\mid l}\ell_we_w=\log l$ over all places above $l$,
Hypothesis~\ref{hyp:twisted-lattices} and the lemma give
\begin{equation}\label{eq:twisted-trim}
  H_j^{\rm cut}\le H_j\le H_j^{\rm cut}+d_j^0\log a_0,\qquad
  H_{j,l}^{\rm cut}\le H_{j,l}.
\end{equation}
The prime divisors of $a_0$ together cost $\log a_0$ once per unit of rank.

\subsection{The two generator primes}

For $e\ge1$ let $\operatorname{Sha}_l(E,e)$ consist of the classes in
$H^1(E,\mathbb Q_l/\mathbb Z_l(-e))$ which vanish on inertia at all
finite places and on the decomposition group at every place above $l$.
Put $t=t_r$. For $l=g_0,g_1$ define, when the lengths are finite,
\[
 V_l=\log l\sum_{e=1}^{t-1}
                 \operatorname{length}_{\mathbb Z_l}\operatorname{Sha}_l(E,e).
\]
The field and the two primes belong to the mesh point. Neither $V_l$
depends on the lift $j$.

\begin{hypothesis}[\openinput]\label{hyp:twisted-generators}
The distinct generator primes satisfy $\log(g_0g_1)\ll y$, as in
\S\ref{sec:finite-correction}. At $l=g_0,g_1$ the groups above are finite.
The original lattices satisfy
the generator-prime comparison of \citep[Hypothesis~17.36]{adriazola-grh}
in the summed form of \citep[\S3]{kimi-generator},
\[
 H_{j,l}\le d_j^0\left[
  \bigl(Cx^C\eta_{\rm sz}+2/l^2+v_l(h_E)\bigr)\log l+V_l\right].
\]
Here $h_E$ is the class number, $C$ is fixed, and the size parameter
satisfies $\log(2+\eta_{\rm sz})\ll_\varepsilon \log x+x^{\varepsilon/4}$.
\end{hypothesis}

\begin{proposition}\label{prop:twisted-generators}
Under Hypothesis~\ref{hyp:twisted-generators}, the trimmed generator
contribution is at most
\[
 H_{j,g_0}^{\rm cut}+H_{j,g_1}^{\rm cut}
       \le d_j^0\bigl(V_{g_0}+V_{g_1}+O_\varepsilon(b)\bigr).
\]
\end{proposition}

\begin{proof}
Apply the surjection of Lemma~\ref{lem:twisted-trim} at each original
place. Since $\log(g_0g_1)\ll y$, their polynomial and size terms total
$o(b)$ since $b=\exp(y)$. The two terms $2\log l/l^2$ are bounded.
The primes are distinct, so their class-number terms sum to at most
$\log h_E$.
Minkowski's bound represents each ideal class by an integral ideal
of norm at most $X=(4/\pi)^{d_E/2}\sqrt{\Delta_E}$. Counting such ideals
gives $h_E\le X^2\zeta_E(2)\le X^2\zeta(2)^{d_E}$.
The degree and $\log\Delta_E$ grow polynomially in $x$ by
\S\ref{sec:heights}, so $\log h_E=o(b)$.
\end{proof}

\subsection{The other coefficient primes}

Fix $l\notin\{g_0,g_1\}$ and a trimmed component. Write $r=k+m-1$,
$n_V=r+t$ and $\epsilon_l$ for the cyclotomic character. A string of
length $s$ in a Galois representation is a sum of the twists
$\sigma\epsilon_l^{-a}$, $0\le a<s$, of one irreducible seed $\sigma$.
The good character string has terms $C_h=\chi\epsilon_l^{-h}$,
$0\le h<t$, for the designated character $\chi$ of the component.

\begin{hypothesis}[\openinput]\label{hyp:twisted-attachments}
For the retained eigensystems the semisimple Hecke--Galois
representations have the string, local inertia and split-prime Hecke
comparisons below \citep{nongenerator-note}.
\begin{enumerate}[label=\textup{(\roman*)},nosep]
\item Group traces and the coefficients of the inertia polynomials
at the level primes belong to the integral Hecke algebra. The
representations are tame there, crystalline with common labelled
weights above $l$, and unramified outside the prescribed exceptional
places and $l$. At the exceptional inertia groups away from $l,g_0,g_1$
the eigencharacters are $1$ ($r$ times) and $\chi$ ($t$ times).
\item Good labels contain the designated string. A retained label
with a global $\chi$-line is good. At a place above $g_0$ the residual
complement of a good string avoids the value of $\chi$ on a tame
generator. In each component meeting the good set the common residual
semisimplification is $\bar\rho\oplus\bigoplus_{h=0}^{t-1}\bar C_h$,
with $\bar\rho$ absolutely irreducible of rank $r>1$.
\item At admissible split places of norm $q$ the integral Hecke
polynomial is the Frobenius polynomial in a common normalization.
After removing the common twist, $C_0$ has modulus $q^{(t-1)/2}$.
The roots of a length-$u$ string are $\alpha q^{(u-1)/2-a}$,
$0\le a<u$, with $q^{-1/2}<|\alpha|<q^{1/2}$ and $|\alpha|=1$
for rank-one seeds. For every finite Galois extension of $\mathbb Q$
containing $E$, one may take all places of $E$ over some rational
prime splitting completely in that extension as test places.
\end{enumerate}
\end{hypothesis}

Assume Hypotheses~\ref{hyp:twisted-lattices}
and~\ref{hyp:twisted-attachments} and the size bound
$\log(g_0g_1)\ll y$ of Hypothesis~\ref{hyp:twisted-generators}.
A summand is potentially abelian if it splits into characters over a
finite extension. Write $m(c)$ for the multiplicity of a residual
character $c$ in the sum of the other summands. A projection costs
at most $a$ if the defect sum in \eqref{eq:twisted-sum} before projection
is at most the sum after projection plus $a$.

\begin{proposition}\label{prop:twisted-strings}
At cost $O_\varepsilon(d_j^0b)$ one may separate the tame inertia
classes and delete all bad full-length strings. On each surviving
bad label of longest length $s<t$,
\[
 I_s=\{0\le h<t:|t-1-2h|\ge s\},\qquad
 m(c)\ge\#\{h\in I_s:\bar C_h=c\}\quad\text{for every }c .
\]
If $s>t/2$, then $t-s$ is even. A surviving label has at most one
rank-one string, of odd length and with centered residual characters.
At $s=t-2$ the long string consists of $C_1,\ldots,C_{t-2}$ in
characteristic zero, and the two outer residual characters occur in
one irreducible summand of rank $2$ or $r+2$.
\end{proposition}

\theoremstyle{plain}
\newtheorem{twistedtest}{Lemma}[theorem]
\renewcommand{\thetwistedtest}{\thetheorem\alph{twistedtest}}

For a number field $L$ of degree $d$ and absolute discriminant
$\Delta$, put
\[
 J(L,n)=(n!)^d\Delta^n
 \left(\prod_{\substack{p\ {\rm prime}\\p\le n+d\ {\rm or}\ p\mid\Delta}}p\right)^n .
\]
This integer clears the denominators of coefficient extraction from unit
values of polynomials of total degree $\le n$ in embedding coordinates
\citep[section on integral coefficient extraction]{nongenerator-note}.
Indeed, for $l>n+d$, $l\nmid\Delta$, a polynomial vanishing on residue
units times the norm polynomial vanishes on the entire residue grid.
Its degree is $<l$, so evaluation has a unit maximal minor. At all other
primes interpolate on $1+l\sum_i j_i\beta_i$, $0\le j_i\le n$, in
an integral basis. Interpolation costs $(n!)^d$ and the inverse
coordinate change costs at most $l^n\Delta^n$.

\begin{twistedtest}[Inertia and exact characters]\label{lem:twisted-exact}
Put $q_v=Nv$, $N_v=\prod_{a=1}^{n_V}(q_v^a-1)$, and
\[
 a_I=\prod_{v\mid g_0g_1}N_v^{n_V},\qquad
 B'=h_EJ(E,2d_E n_V).
\]
At cost $d_j^0\log\bigl(a_I\,t!(B')^t\bigr)$ one may split the
inertia classes and require every line of a rank-one string of length
$u>r$ to equal some $C_h$.
\end{twistedtest}

\begin{proof}
Frobenius permutes the tame roots by $z\mapsto z^{q_v}$, so they lie
in $\mu_{N_v}$. For $\xi,z\in\mu_{N_v}$,
\[
 \sum_{i=0}^{N_v-1}(z/\xi)^i=N_v\,\mathbf1_{z=\xi}.
\]
Multiply over the roots, sum over distinct orderings of each prescribed
multiset, and multiply over $v$. The resulting symmetric polynomials
in the inertia roots give $a_I\pi_H$ integrally for every class $H$
after adjoining the roots of unity. The length normalization descends.
For $M'=\bigoplus_H\pi_HM$ one has $a_I M'\subset M\subset M'$.
The indices on the good projection and on the good intersection
each have length at most $v(a_I)\dim V_G$. Their difference gives the
change in defect.

In a class meeting $G$, let $\tau$ be the chosen tame generator over
$g_0$. The factor $(X-\chi(\tau))^t$ is coprime over $O$ to its
complement, giving a rank-$t$ projector $P$. Each line $\nu$ in a
rank-one string of length $u>r$ lies in its image. At every exceptional
inertia group the common value on these $u$ lines must equal $\chi$,
since at most $r$ roots differ. Thus $\lambda=\nu/C_h$ is unramified
away from $l$. The intervals of $u$ and $t$ labelled weights overlap
in a multiset of size $n_V$, so its exponents satisfy $|k_\tau|\le n_V$.

On local units $\lambda(v)=\prod_\tau\tau(v)^{k_\tau}$
\citep[Proposition~B.4]{conrad-local-characters}. In
$\mathrm N(v)^{n_V}\lambda(v)$ extract the monomial
$\prod_\tau\tau(v)^{n_V}$ with denominator $J(E,2d_E n_V)$.
Only everywhere unramified ratios remain. Summing their values over
$\operatorname{Cl}(E)$ keeps the trivial ratio. Twisting each unit
and class representative by $C_h^{-1}$ thus gives $D_h\in O[G_E]$
with
\[
 D_h|_\nu=B'\,\mathbf1_{\nu=C_h}
 \qquad(\nu\text{ a good line or a line as above}).
\]
On a good label $P(\sum_h D_h)P$ equals $B'P$. A line unmatched
by all $C_h$ lowers its rank below $t$. By Newton's identities,
$t!$ times its $t$-th elementary symmetric function is an integral
Hecke operator. Its good value is $t!(B')^t$. Apply
Lemma~\ref{lem:twisted-trim}.
\end{proof}

\begin{twistedtest}[Position of a character]\label{lem:twisted-position}
Put $B=J(\mathbb Q,4d_Et)$. There are $Z_h\in O[G_E]$ with
$Z_h|_{C_a}=B\,\mathbf1_{a=h}$. On a potentially abelian irreducible
summand of a short label, each $Z_h$ acts as $0$ or $B$.
At most one position gives $B$, and none in $I_s$ does.
\end{twistedtest}

\begin{proof}
Choose a common Galois splitting field $L/\mathbb Q$, unramified over
$E$ above $l$. Indeed, if $U$ splits into lines over $L'/E$ Galois,
the finite conjugation action on $\operatorname{End}_{G_{L'}}U$ is
crystalline above $l$ and hence unramified there. Its kernel acts by
scalars on each isotypic space. Take the compositum of the resulting
fields and their conjugates.

For a constituent line $\nu$ over $L$, set $\lambda=\nu/C_h|_{G_L}$.
The unit theorem makes the sum of conjugate labelled exponents of
$\lambda$ a constant $W$. At a completely split test prime $p$,
reciprocity gives
\[
 p^{-[L:\mathbb Q]W/2}
    =\prod_{v\mid p}|\lambda(\operatorname{Frob}_v)|,\qquad |W|<2t .
\]
Indeed all normalized short roots have modulus in
$(p^{-s/2},p^{s/2})$, and $C_h$ has modulus $p^{(t-1-2h)/2}$.
For $h\in I_s$ every ratio is on the same side of $1$, hence $W\ne0$.

In the fixed field of complex conjugation the unit character is the
norm to the power $W$. At a place above $l$ of ramification index
$e\le d_E$, choose a unit of norm $z\in\mathbb Z_l^\times$ in its
maximal unramified subfield. Its full norm is $z^e$. Lift the
corresponding Artin element to $g_z\in G_L$ and put $A=2d_Et$.
Coefficient extraction from the group-algebra samples
$z^AC_h(g_z)^{-1}[g_z]$ acts on $\nu$ by
\[
 B[z^A]z^{A+eW}=B\,\mathbf1_{W=0}.
\]
This gives $Z_h$. Conjugate lines of an irreducible summand have the
same $W$. The good weights are spaced by two, so at most one $h$
matches. On $C_a$ this occurs precisely at $h=a$.
\end{proof}

\begin{twistedtest}[Residual multiplicities]\label{lem:twisted-multiplicities}
Put $I_c=\{h:\bar C_h=c\}$ and $D=t!B^t$. At cost
$d_j^0\log D^{2t}$ every one-dimensional residual factor of a
potentially abelian summand occupies a distinct matching position
of its own type. This includes $l=2$.
\end{twistedtest}

\begin{proof}
Lift the idempotent for $c$ in the semisimple residual group-algebra
image. Hensel factorization of the product of the characteristic
polynomials of the lift gives an integral projector $P_c$ onto the
$1$-primary part on every label. Its rank on each summand equals
the multiplicity of $c$ there. For $\varnothing\ne I\subset I_c$ put
\[
 k=|I|,\qquad U_I=P_c\Big(\sum_{h\in I}Z_h\Big)P_c,\qquad
 T_I=k!\operatorname{tr}\bigl(\textstyle\bigwedge^k U_I\bigr).
\]
Newton's identities give $T_I$ in the integral Hecke algebra. Its
good value is $k!B^k$. It vanishes whenever $\operatorname{rank}U_I<k$.

Take $I=I_c$. An unmatched copy of $c$ gives rank $<|I_c|$.
Next omit $a\in I_c$. If two copies match $a$, the rank is at most
$|I_c|-2$ (for $|I_c|=1$ there is only one copy to begin with).
There are at most $2t$ tests, and every $k!B^k$ divides $D$.
Multiply the rescaled tests and pad with constant factors to good
value $D^{2t}$, then apply Lemma~\ref{lem:twisted-trim}.
For the endpoint bound alone the sets $I_c\cap I_s$ give at most
$t^2$ tests and good value $D^{t^2}$.
\end{proof}

\begin{proof}[Proof of Proposition~\ref{prop:twisted-strings}]
The three lemmas cost $O_\varepsilon(d_j^0b)$ since
$\log a_I\ll n_V^3d_E\log(g_0g_1)$ and the logarithms of $B',B,t!$
grow polynomially in $x$. No potentially abelian factor occupies a
position in $I_s$, which gives the bound for $m(c)$ with multiplicity.

A string of length $s\ge2$ repeats each residual seed constituent
$s$ times up to twist. Only one constituent has rank $>1$. Thus
$s\dim\sigma\le t$, and $s>t/2$ forces rank one. For a length-$t$
string Lemma~\ref{lem:twisted-exact} gives $t$ distinct good lines.
One is $\chi$, so the label is good. For a short length $s>t/2$,
\[
 s+|I_s|=s+2\lceil(t-s)/2\rceil\le t .
\]
Hence $t-s$ is even. In particular $s=t-1$ is excluded.

In a rank-one string of any length $u$, line $a$ has modulus
$p^{(u-1)/2-a}$, so its position is $h=(t-u)/2+a$.
As $t$ is odd, $u$ is odd and the middle position is occupied.
Lemma~\ref{lem:twisted-multiplicities} allows it only once.

At $s=t-2>r$ the exact lines are therefore $C_1,\ldots,C_{t-2}$.
The two outer factors lie in non-potentially-abelian summands.
If they lay in separate summands each would need the unique
$\bar\rho$ to have rank greater than one. Their common summand has
rank $2$ or $r+2$. Here $t\ge r+5$ and $t>4$ by \S\ref{sec:theta-chain}.
\end{proof}

A short label whose distinguished $t$-dimensional inertia space is
a global sum of character lines has $m(c)=0$ for every $c$.
The endpoint bound deletes it even if $\bar C_0=\bar C_{t-1}$.

\begin{corollary}\label{cor:twisted-ray}
In the length-$(t-2)$ case of Proposition~\ref{prop:twisted-strings},
suppose $l>2$ is unramified in $E$ and $\bar C_0\ne\bar C_{t-1}$.
At every place over $l$ suppose a common crystalline twist puts the
good and bad weights in one Fontaine--Laffaille interval of width
at most $l-2$. Let $S$ be a finite set of primes of $E$, disjoint
from those above $l$, containing every other allowed ramified prime.
Put $E_1=E(\mu_l)$, $\Gamma_1=\operatorname{Gal}(E_1/E)$, and let
$\mathfrak m$ be the product of the primes of $E_1$ above $S$.
If the summand carrying the endpoints has rank two, then, for the
splitting residue field $\mathbb F$,
\[
 \operatorname{Hom}_{\mathbb F_l[\Gamma_1]}\!\left(
 \frac{\operatorname{Cl}_{\mathfrak m}(E_1)}
 {l\operatorname{Cl}_{\mathfrak m}(E_1)+\langle[v]:v\mid l\rangle},
 \mathbb F(\bar\epsilon_l^{\,1-t})\right)\ne0 .
\]
\end{corollary}

\begin{proof}
The rank-two lattice lemma gives a stable lattice whose reduction
is a nonsplit extension of $\bar C_0$ by $\bar C_{t-1}$
\citep[Theorem~1]{bellaiche-ribet}. Its coefficient character is
$\bar C_{t-1}/\bar C_0=\bar\epsilon_l^{\,1-t}$.
In the covariant Fontaine--Laffaille convention
$\epsilon_l$ has weight $-1$. After the common twist, the filtration
jumps of the subline and quotient at $\tau$ are $h_\tau+t-1$ and
$h_\tau$. The exact, fully faithful torsion functor identifies
these lines with the reductions of the good characters
\citep[Fact~4.10 and Remark~4.11]{booher-fl}. In the torsion module
$\varphi^{h_\tau}$ kills the higher subline. Its image on any lift of
the lower line maps isomorphically to the quotient at the next
embedding. These image lines form a filtered Frobenius-stable
complement, so the extension splits at every place over $l$.

Since $|\Gamma_1|$ is prime to $l$, restriction to $E_1$ preserves
the nonzero class. There it is an additive character, trivial at
the places over $l$. At each prime over $S$ it kills the principal
units, whose residue characteristic differs from $l$. Artin
reciprocity gives the indicated homomorphism. Vanishing of this
space would exclude the rank-two summand at that component.
\end{proof}

\subsection{The remaining integral matrices}

After the deletions, let $M'$ be the projected lattice of a mixed
component and write $KM'=V_G\oplus V_B$ for its good and surviving
bad spaces. Choose an admissible split unramified test place $v_0$.
Evaluate the integral Hecke polynomial there at
$\chi(\operatorname{Frob}_{v_0})$ in the matching normalization,
and denote the endomorphism of $M'$ by $P$. Each good label has this
root. Every surviving bad label has longest string of length $s<t$,
and the modulus of every normalized root is strictly smaller than
the modulus $q^{(t-1)/2}$ of the good endpoint. Thus
$\ker_K P=V_G$ and $\operatorname{im}_K P=V_B$.

\begin{lemma}\label{lem:twisted-smith}
Let $d_B=\dim_K V_B$ and $I_{d_B}(P)$ be the ideal of $d_B$-rowed
minors of $P$ on $M'$. Then
\begin{equation}\label{eq:twisted-smith}
 e_O\eta(M',G)=v(\det(P|_{V_B}))-v(I_{d_B}(P)).
\end{equation}
For $d_B=0$ take the determinant and the ideal generator to be $1$.
\end{lemma}

\begin{proof}
Write $L_B=M'\cap V_B$. Projection on either space identifies the
defect module with $M'/((M'\cap V_G)+L_B)$. On the bad space it
therefore equals $\pi_B M'/L_B\simeq PM'/P L_B$. In
$P L_B\subseteq PM'\subseteq L_B$ the outer index has length
$v(\det(P|_{V_B}))$. The saturation of $PM'$ in $M'$ is $L_B$, so
Smith normal form gives $\operatorname{length}_O(L_B/PM')=v(I_{d_B}(P))$.
Subtracting proves \eqref{eq:twisted-smith}.
\end{proof}

For each lift let $R_j^{\rm tw}$ be the sum of the right side of
\eqref{eq:twisted-smith} over the surviving mixed components at
$l\ne g_0,g_1$, with weight $\ell_we_w/e_O$ for each. Pure
components have zero defect. At each component one split test place
suffices. The difference of valuations computes the length for any
such choice. In particular $R_j^{\rm tw}$ includes the surviving
components at $2$. The ray class test requires its additional local
hypotheses, and its vanishing on the selected family is unknown.

\begin{hypothesis}[\openinput]\label{hyp:twisted-residual}
On the chosen mesh points there are constants independent of the
window, point, shift and lift, after fixing $\varepsilon$, such that
\[
 V_{g_0}+V_{g_1}\ll_\varepsilon y(B_0+N_c),\qquad
 R_j^{\rm tw}\ll_\varepsilon d_j^0\, y(B_0+N_c)\quad(1\le j\le m).
\]
\end{hypothesis}

\begin{theorem}\label{thm:twisted-defect}
Under Hypotheses~\ref{hyp:twisted-lattices},
\ref{hyp:twisted-generators} and~\ref{hyp:twisted-attachments}, at
the chosen shift one has
\[
 H_j\le R_j^{\rm tw}
          +d_j^0\bigl(V_{g_0}+V_{g_1}+O_\varepsilon(b)\bigr)
                       \qquad(1\le j\le m).
\]
Under Hypothesis~\ref{hyp:twisted-residual} it follows that
$H_j\ll_\varepsilon d_j^0\,y(B_0+N_c)$ uniformly on the windows.
\end{theorem}

\begin{proof}
Start with \eqref{eq:twisted-trim}. Proposition~\ref{prop:twisted-generators}
bounds the two generator terms on the original support.
Proposition~\ref{prop:twisted-strings} bounds the cost of passing to
the other projected lattices, and Lemma~\ref{lem:twisted-smith}
computes their remaining lengths. The term $\log a_0$ and every
logarithmic deletion cost are $O_\varepsilon(b)$ since $b=\exp(x^\varepsilon)$.
Finally $b\le y(B_0+N_c)$.
\end{proof}

For the target of \S\ref{sec:heights} let $D_0>0$ be the number of
base rows. If its frame has rank $D_0$ and the level-$j$ row spaces
have dimension $D_0$, then $d_j^0\le D_0$ by the adjoint image
definition. The base determinant comparison divides the local index
sum by $D_0$ \citep[Proposition~15.53]{adriazola-grh}. Suppose the
twisted indices of this frame are bounded by a uniform multiple of
these defects. Under the four hypotheses of this section their
contribution is then bounded using
\[
 D_0^{-1}\sum_{j=1}^m H_j
    \ll_\varepsilon m y(B_0+N_c)\le y\mathcal B .
\]
This is the twisted part of the target-slope allowance in
Hypothesis~\ref{hyp:height-norms}. The period comparison of
Hypothesis~\ref{hyp:period-frame} uses these same lattices and the
rank-weighted heights of the words. Both frame identifications still
require proof on the selected family. The finite correction $E_u$
of \S\ref{sec:finite-correction} already has its own place sum
\eqref{eq:finite-place-sum}.
